Integral closure in a field commutes with localization. Indeed, if an element is integral over $A_s$, multiplying it by a sufficiently high power of $s$ gives an element integral over $A$; the converse follows by localizing a monic equation.

Thus, if $\overline A$ denotes the integral closure of $A$ in $K(X)$, the inverse image of $D(s)\subseteq\operatorname{Spec}A$ in $\operatorname{Spec}\overline A$ is canonically $\operatorname{Spec}\overline{A_s}$. Using common distinguished neighborhoods from \exref{II.3.1}, these identifications glue the affine normalizations along every overlap. They satisfy the cocycle condition because all ring maps come from inclusions in $K(X)$. The resulting scheme $\widetilde X$ is integral and normal, and the affine maps $\operatorname{Spec}\overline A\to\operatorname{Spec}A$ glue to $\widetilde X\to X$.

For a dominant morphism $f:Z\to X$ with $Z$ normal and integral, the induced inclusion $K(X)\to K(Z)$ sends $\overline A$ into every affine coordinate ring of $Z$ lying over $\operatorname{Spec}A$: these rings are integrally closed, and the images of elements of $\overline A$ are integral over them. Hence $f$ factors locally through the affine normalizations. These factorizations agree on overlaps and are unique because they induce the fixed inclusion of function fields. They therefore glue uniquely.

If $X$ is of finite type over $k$, every such $A$ is a finitely generated $k$-domain, whose integral closure is a finite $A$-module. The normalization morphism is consequently finite.
